\documentclass[11pt]{article}
\usepackage[margin=1.15in]{geometry}
\usepackage{amsmath,amssymb,amsthm,mathtools}
\usepackage{xcolor}
\usepackage[colorlinks=true,linkcolor=blue!60!black,citecolor=blue!60!black,urlcolor=blue!60!black]{hyperref}

\newtheorem{theorem}{Theorem}[section]
\newtheorem{proposition}[theorem]{Proposition}
\newtheorem{lemma}[theorem]{Lemma}
\newtheorem{corollary}[theorem]{Corollary}
\newtheorem{cproposition}[theorem]{Computational Proposition}
\theoremstyle{definition}
\newtheorem{definition}[theorem]{Definition}
\theoremstyle{remark}
\newtheorem{remark}[theorem]{Remark}

\newcommand{\C}{\mathbb{C}}
\newcommand{\R}{\mathbb{R}}
\newcommand{\dg}{\dagger}
\newcommand{\dd}{\ddagger}
\newcommand{\ad}{\operatorname{ad}}
\newcommand{\ob}{\mathfrak o}
\newcommand{\Tr}{\operatorname{Tr}}
\newcommand{\Res}{\operatorname*{Res}}
\newcommand{\AK}{\mathcal A_K}
\newcommand{\Ap}{\mathcal A_+}
\newcommand{\PiE}{\Pi_{\ker\ad_{h_0}}}

\title{Interaction Obstructions in Einstein--Weyl Gravity:\\
Cubic Protection, Second-Order Metric Failure, and Krein Visibility}
\author{GPT-5.6.sol\thanks{OpenAI model.  Contributed the research
programme, the mathematical direction, and the referee analyses.}
\and
Claude Fable 5\thanks{Anthropic model.  Contributed the derivations, the
machine verification (SymPy and exact rational arithmetic), and the
manuscripts.  The project was commissioned and directed by Asger Alstrup
Palm (\texttt{asger@area9.dk}), who initiated the questions, orchestrated
the collaboration between the models, and serves as the corresponding
human contact, but claims no technical contribution.}}
\date{}

\begin{document}
\maketitle

\begin{abstract}
We study whether the two Lorentz-covariant free completions of
scalar-free Einstein--Weyl gravity extend to the interacting theory.  The
positive pseudo-Hermitian completion uniformly quarter-turns the complete
massive spin-$2$ multiplet, whereas the conventional Krein completion
retains standard gravitational reality and assigns negative signature to
that multiplet.  We first prove a cubic-protection theorem.  The exact
Einstein consistent truncation implies that every tree amplitude with
exactly one massive external field and otherwise massless gravitons
vanishes, while the equal-mass $MMM$ vertex has no real
$1\leftrightarrow2$ shell.  Hence the first-order positive-metric
obstruction is zero.

At second order the physical process $MM\to Mh$ is open.  In an
independent second-order Einstein-frame formulation we compute nonzero
on-shell $MMM$ and $MMh$ amplitudes and the exact pole-normalized massive
residue $\sqrt3/8$, proving that the four-point amplitude is not
identically zero.  A complete calculation in the original fourth-order
theory --- including the quartic contact term, all $s,t,u$ exchanges, the
physical adjoint, and gauge and equation-of-motion checks --- gives at an
interior rational physical point
\begin{equation*}
\PiE\bigl(T_{2,+}^{\dg}-T_{2,+}\bigr)
=-\frac{15762482064\,i}{5584765625}\,\sigma_x\ne0 .
\end{equation*}
The unique Lorentz-covariant positive free metric therefore admits no
analytic second-order deformation compatible with the interaction.

We then classify conventional Poincar\'e-covariant, BRST-compatible,
cluster-factorizing Krein gradings.  Agreement with the free gravitational
signature fixes their Fock lift to $J_F=(-1)^{N_M}$.  The gravitational
obstruction survives BRST cohomology and is not null:
\begin{equation*}
\mathcal O^\sharp=\mathcal O,
\qquad
\Tr(\mathcal O^\sharp\mathcal O)=-8\AK^2\ne0 .
\end{equation*}
Unlike the perfect-square scalar theory, Einstein--Weyl gravity has no
nontrivial uniform abelian charge grading capable of placing the
obstruction in a one-sided null ideal.  The result excludes both an
analytic continuation of the canonical positive completion and the
conventional massive-number grading as a null-sector probability
mechanism.  It does not exclude indefinite Krein evolution itself, nor
nonlocal, non-factorizing, nonperturbative, or conformal-boundary
completions.
\end{abstract}

\tableofcontents

\section{Introduction}\label{sec:intro}

Scalar-free Einstein--Weyl gravity is the simplest four-dimensional
quadratic-curvature theory whose flat-space physical spectrum consists of
the ordinary massless graviton and one additional massive spin-$2$
multiplet.  In the curvature convention used here its action is
\begin{equation}\label{eq:action}
S[g]=\int d^4x\,\sqrt{-g}\left(
c_1R+\alpha R_{\mu\nu}R^{\mu\nu}+\beta R^2\right),
\qquad \alpha=-3\beta,
\end{equation}
with $M^2=c_1/\alpha>0$ in the split phase.  Up to the four-dimensional
Euler density, the quadratic term is Weyl-squared.  The extra spin-$2$
branch has the opposite kinetic sign, the familiar higher-derivative
ghost of quadratic gravity \cite{Stelle}.

Paper~IV of this series \cite{paperIV} did not choose between proposed
interacting cures.  It classified the \emph{free} covariant completions.
After symplectic reduction, all five massive polarizations form one
irreducible massive spin-$2$ representation.  Covariance therefore
allows no helicity-by-helicity compromise.  There are two real-form
classes in the stated free, quasifree, mode-local setting:
\begin{enumerate}
\item a positive pseudo-Hermitian completion in which the entire massive
multiplet has quarter-turned reality; and
\item a Krein completion with standard gravitational reality, positive
massless signature, and uniform negative massive signature.
\end{enumerate}
They determine the same free complex spectral functional but different
adjoints and probability structures.  The natural interacting question
is therefore sharp:
\begin{quote}
Does either covariant free completion remain compatible with the full
Einstein--Weyl interaction?
\end{quote}

Paper~V \cite{paperV} supplied the deformation complex needed to ask the
positive part of this question.  If $\rho_0$ is the free Dyson map and
$v_n$ are the transported interaction vertices, an analytic correction
$\rho_\zeta=\exp[-\tfrac12\sum_n\zeta^nR_n]\rho_0$ must solve
\begin{align}\label{eq:defcomplex}
[h_0,R_1]&=v_1^{\dg}-v_1,\nonumber\\
[h_0,R_2]&=(v_2^{\dg}-v_2)
+\frac12[R_1,v_1+v_1^{\dg}].
\end{align}
On a degenerate free-energy shell the right-hand side must have zero
projection onto $\ker\ad_{h_0}$.  That shell projection is the metric
obstruction.  In the perfect-square scalar, first-order protection was
followed by a second-order branch-changing obstruction; at the exact
massless $O(1,1)$ boundary, however, the offending component could become
one-sided in continuous charge and null in the generalized Born trace.

Gravity has the same perturbative-order pattern and a different
mechanism:
\begin{equation}\label{eq:pattern}
\boxed{\ob_+^{(1)}=0,\qquad \ob_+^{(2)}\ne0.}
\end{equation}
The cubic protection follows from the exact Einstein subsector together
with equal-mass kinematics.  The second-order obstruction is carried by
the physically open process
\begin{equation}\label{eq:process}
M+M\longrightarrow M+h,
\end{equation}
which changes massive particle number by one.  The same process also
separates gravity from the scalar boundary mechanism.  Covariance and
cluster factorization fix the conventional gravitational fundamental
symmetry to $(-1)^{N_M}$, but $\mathbb Z_2$-odd does not imply null:
the product of two odd operators is neutral and its quadratic trace is
nonzero.

\paragraph{Results.}
The proof has four independent layers.
\begin{enumerate}
\item Section~\ref{sec:phase} proves a diagram-level phase theorem: a
connected amplitude with $E_M$ external massive legs changes by
$(-i)^{E_M}$ under the uniform quarter-turn, including the compensating
sign of every internal massive propagator.
\item Section~\ref{sec:cubic} proves cubic protection.  The Einstein
consistent truncation kills $Mhh$, while $MMM$ has no real physical
$1\leftrightarrow2$ shell.
\item Sections~\ref{sec:factorization}--\ref{sec:full} prove the
second-order obstruction twice: first by an independent nonzero
Einstein-frame factorization residue, then by a complete exact physical
four-point calculation in the original fourth-order theory.
\item Sections~\ref{sec:grading}--\ref{sec:brst} prove conventional
Krein visibility: the covariant Fock grading is unique in the stated
class, the exact block is not null, no scalar-like abelian charge exists,
and the block is not BRST-exact.
\end{enumerate}

\paragraph{Comparison with the scalar obstruction.}
The parallel and the difference are summarized by
\begin{center}
\begin{tabular}{p{0.42\linewidth}|p{0.42\linewidth}}
\hline
Perfect-square scalar & Einstein--Weyl gravity\\
\hline
K\"all\'en zero and reality-factor protection & Einstein consistent
truncation and equal-mass kinematics\\
$H+L\to L+L$ & $M+M\to M+h$\\
$O(1,1)$ one-sided null ideal at the exact boundary & no nontrivial
uniform abelian charge grading\\
obstruction can become probability-invisible & obstruction remains
non-null for the conventional covariant Krein form\\
\hline
\end{tabular}
\end{center}

\paragraph{Strength of the claims.}
``Metric failure'' means failure of the analytic order-by-order
deformation of the canonical covariant positive free metric in the
operator class of \eqref{eq:defcomplex}.  It is not an absolute theorem
against every positive metric, and no complex interacting spectrum is
proved.  ``Krein visibility'' has the precise finite-block meaning
$\mathcal O^\sharp\mathcal O\ne0$ and
$\Tr(\mathcal O^\sharp\mathcal O)\ne0$ for the unique conventional
covariant, cluster-factorizing grading.  It is a non-nullity statement,
not a claim that an indefinite trace is a positive probability.  In
particular, the paper does not rule out an indefinite Krein dynamics;
it rules out the naive parity/null-sector mechanism as a replacement for
a positive interacting probability calculus.

\paragraph{Verification.}
Checks G13--G18 are exact SymPy or exact rational computations.  They are
mapped to the analytic statements in Appendix~\ref{app:checks}.  The
exhaustive $5^3\times2$ polarization scan proposed as G16 is not used:
the factorization residue, exact real-shell certificate, crossing, Ward,
and gauge checks already close the proof.  The scan remains optional
regression hardening.

\section{Free input, conventions, and the deformation class}
\label{sec:free}

\subsection{Physical one-particle spectrum}

We use signature $(+,-,-,-)$ in the main text.  Around flat space the
reduced physical one-particle space is
\begin{equation}\label{eq:oneparticle}
\mathcal H_{\rm phys}^{(1)}=
\mathcal H_{h,+2}\oplus\mathcal H_{h,-2}\oplus\mathcal H_M,
\qquad \dim\mathcal H_M=5.
\end{equation}
The massless summands carry helicities $\pm2$ on $p^2=0$; the massive
summand carries the spin-$2$ representation on $p^2=M^2$.  Paper~IV
derives this by explicit diffeomorphism reduction rather than by a
covariant propagator count.  Its result is a healthy massless sector and
a uniformly ghost-signed massive sector.

\subsection{Two free real forms}

In the positive class the original massive mass eigenfield is
$\dd$-anti-Hermitian.  We write the positive-frame observable as
\begin{equation}\label{eq:quarterturn}
\widehat M_{\mu\nu}=iM_{\mu\nu},
\qquad M_{\mu\nu}=-i\widehat M_{\mu\nu}.
\end{equation}
The same quarter-turn acts on every one of the five polarizations.  The
massless graviton retains standard reality.  In the conventional Krein
class both fields retain standard gravitational reality, while the
one-particle fundamental symmetry has signature
\begin{equation}\label{eq:freesignature}
J^{(1)}_{\rm free}=I_2\oplus(-I_5),
\qquad \text{signature }(+,+;-,-,-,-,-).
\end{equation}

The uniformity is not optional.  Since $\mathcal H_M$ is an irreducible
massive spin-$2$ representation, Schur's lemma makes every covariant
Hermitian form on its fiber a scalar multiple of the identity.  Rotating
only selected helicities would violate Lorentz covariance
\cite{paperIV}.

\subsection{Finite-volume shell projection}

As in Paper~V, the cohomological equation is interpreted first at finite
volume, where momenta are discrete and $h_0$ has point spectrum.  If
$P_E$ projects onto free energy $E$, then
\begin{equation}\label{eq:kernelproj}
\PiE S=\sum_E P_ESP_E,
\end{equation}
and $[h_0,R]=S$ has a solution only if this projection vanishes.  The
rational physical point below can be placed on a momentum lattice after
an overall choice of box scale.  Continuum statements mean that the
corresponding on-shell matrix element is nonzero on an open subset of the
real scattering shell; a distributional rigged-Fock reformulation is not
needed for the present existence theorem.

\subsection{Amplitude convention}

The perturbiner returns the multilinear coefficient of the action with
the universal Feynman $i$ stripped.  We call it the reduced tree
amplitude.  At real momenta and real linear polarizations it is real.  A
global action rescaling and external polarization rescalings change the
displayed rational certificate but cannot change its nonvanishing.  All
statements involving the precise coefficient below are therefore made in
the conventions of Appendix~\ref{app:fourpoint}; the obstruction theorem
itself is the invariant nonzero statement.

\section{The external-massive-leg phase theorem}\label{sec:phase}

The quarter-turn phase must be transported through complete diagrams,
not appended to an already computed external state by slogan.  Internal
massive lines carry the compensating kinetic sign.

\begin{theorem}[External-massive-leg phase]\label{thm:phase}
Let $\mathcal D$ be a connected perturbative diagram written in
mass-eigenfields $(h,M)$, with $E_M$ external massive legs and $I_M$
internal massive lines.  Under \eqref{eq:quarterturn}, its complete
reduced amplitude obeys
\begin{equation}\label{eq:phasetheorem}
\mathcal A_+^{(E_M)}(\mathcal D)
=(-i)^{E_M}\mathcal A_K^{(E_M)}(\mathcal D).
\end{equation}
The statement includes contact terms, exchanges, and the sign of each
massive inverse kernel.
\end{theorem}

\begin{proof}
Let $n_v$ be the number of massive fields incident at vertex $v$.  Field
substitution contributes
\begin{equation}
\prod_v(-i)^{n_v}.
\end{equation}
The massive quadratic form changes sign because
$(-i)^2=-1$; hence its inverse, the internal massive propagator, also
contributes $-1$ per internal massive line.  Every internal massive line
is incident twice, so
\begin{equation}
\sum_v n_v=E_M+2I_M.
\end{equation}
The total factor is therefore
\begin{equation}
(-i)^{E_M+2I_M}(-1)^{I_M}
=(-i)^{E_M}(-1)^{I_M}(-1)^{I_M}
=(-i)^{E_M}.
\end{equation}
Massless propagators are unchanged.  Since the argument is diagramwise,
it applies to their complete sum.  Local invertible changes between
regular mass-eigenfield frames can redistribute contact and exchange
terms but leave the on-shell statement unchanged.
\end{proof}

\begin{corollary}[Phase selection]\label{cor:phase}
For a real standard-frame tree amplitude, even $E_M$ gives a real
positive-frame coefficient and odd $E_M$ gives a purely imaginary one.
Only odd-$E_M$ processes can contribute to the ordinary-adjoint
anti-Hermitian shell source.
\end{corollary}

For the four-point process \eqref{eq:process}, $E_M=3$, and
\begin{equation}\label{eq:mmmhphase}
\mathcal A_+=i\mathcal A_K.
\end{equation}
The sign is tied to the convention $M=-i\widehat M$; reversing the
quarter-turn reverses the displayed sign but not the obstruction's
nonvanishing.

\section{Cubic protection}\label{sec:cubic}

There are two odd-massive cubic structures: $Mhh$ and $MMM$.  The first
is eliminated dynamically, the second kinematically.

\subsection{Einstein consistent truncation}

The scalar-free equations may be written schematically as
\begin{equation}\label{eq:ewfield}
c_1G_{\mu\nu}+\gamma_B B_{\mu\nu}=0,
\end{equation}
where $B_{\mu\nu}$ is the Bach tensor and $\gamma_B$ depends on the
normalization in \eqref{eq:action}.  Every Ricci-flat Einstein metric has
$G_{\mu\nu}=0$ and is Bach-flat.  More generally Einstein metrics are
Bach-flat \cite{LiuLuPope}.  Thus the nonlinear Einstein solution space
is an exact subsector of the Einstein--Weyl equations.

On the regular split branch, the higher-derivative theory admits a local
second-order mass-eigenfield formulation in which the Einstein metric and
the nonlinear massive spin-$2$ field are separate variables
\cite{MagnanoSokolowski}.  Denote them by $(h,M)$ near flat space.  The
field map is algebraic and locally invertible on this branch, and $M=0$
is precisely the Einstein subsector.  These hypotheses are what make the
following LSZ statement stronger than a basis-dependent observation.

\begin{lemma}[One-massive-leg rule]\label{lem:onemassive}
For asymptotic mass eigenfields on the regular split branch,
\begin{equation}\label{eq:onemassive}
\mathcal A_{\rm tree}(M,h,\ldots,h)=0.
\end{equation}
In particular $\mathcal A_3(Mhh)=0$ for every physical polarization.
\end{lemma}

\begin{proof}
Write the massive-field equation in the second-order frame as
\begin{equation}
\mathcal K_M M+J_M^{(2)}[h,h]+J_M^{(3)}[h,h,h]+\cdots=0.
\end{equation}
For every nonlinear vacuum Einstein solution $h_{\rm E}$, the pair
$(h_{\rm E},M=0)$ solves the full theory.  Hence the terms independent of
$M$ vanish on the Einstein graviton shell.  Functional differentiation
and LSZ reduction say that a physical amplitude with one external $M$
and all remaining external legs in the Einstein subsector is zero.
Terms generated by a local invertible change of fields are proportional
to equations of motion or external inverse propagators and vanish under
LSZ, the standard equivalence-theorem mechanism \cite{KOS}.
\end{proof}

The exact perturbiner check evaluates the physical decay point with the
massive particle at rest and two back-to-back massless gravitons.  All
$5\times4=20$ combinations of the five massive polarizations and four
graviton-helicity choices vanish exactly.  This is a check of the
amplitude-level content of Lemma~\ref{lem:onemassive}, not its only
justification.

\subsection{Equal-mass cubic kinematics}

The $MMM$ vertex itself is nonzero at complex on-shell kinematics
(Section~\ref{sec:factorization}).  It nevertheless has no real physical
$1\leftrightarrow2$ shell.  If $p_3=p_1+p_2$ with future-directed
$p_i^2=M^2$, then
\begin{equation}
p_3^2=(p_1+p_2)^2\ge(2M)^2,
\end{equation}
which cannot equal $M^2$.  Equivalently, a particle of mass $M$ cannot
decay into two particles of the same mass.

\begin{theorem}[Cubic protection]\label{thm:cubicprotection}
For scalar-free Einstein--Weyl gravity about flat space, the
Lorentz-covariant positive real form has no first-order physical on-shell
metric obstruction:
\begin{equation}\label{eq:firstzero}
\boxed{\ob_+^{(1)}=0.}
\end{equation}
\end{theorem}

\begin{proof}
By Corollary~\ref{cor:phase}, only odd-massive cubic amplitudes can enter
$\PiE(v_1^{\dg}-v_1)$.  The $Mhh$ amplitude vanishes by
Lemma~\ref{lem:onemassive}; the $MMM$ amplitude has no real physical
shell.  The remaining cubic structures carry an even number of massive
legs.  Therefore the first-order anti-Hermitian source has zero physical
shell projection.
\end{proof}

This protection is structurally different from the K\"all\'en-factor
zero of Paper~V.  It is a consequence of an exact nonlinear consistent
truncation plus massive-particle kinematics.

\section{Independent Einstein-frame factorization rail}
\label{sec:factorization}

The next odd-massive process is four-point $MM\to Mh$.  Before expanding
the original fourth-order action to quartic order, we prove independently
that this amplitude is not identically zero.

\subsection{Second-order action and convention change}

This section follows the exact Einstein-frame formulation of Magnano and
Soko\l{}owski \cite{MagnanoSokolowski}.  Its verification artifact uses
signature $(-,+,+,+)$, only in this section and Appendix~\ref{app:eframe};
thus a massive pole is $P^2=-M^2$.  With $M=1$ and in the notation of that
formulation, the nonlinear massive-field Lagrangian is
\begin{equation}\label{eq:LM}
L_M=\frac12K(Q)+\frac{m^2}{8A}
\left(-\tau^2+\psi_{\mu\nu}\psi^{\mu\nu}
+6A\tau-12A^2\right),
\qquad \psi=g+\Phi .
\end{equation}
Here $Q^\alpha{}_{\mu\nu}=\Gamma^\alpha{}_{\mu\nu}(\psi)
-\Gamma^\alpha{}_{\mu\nu}(g)$,
\begin{equation}
K(Q)=g^{\mu\nu}\left(
Q^\alpha{}_{\mu\nu}Q^\beta{}_{\alpha\beta}
-Q^\alpha{}_{\mu\beta}Q^\beta{}_{\nu\alpha}\right),
\end{equation}
$A$ is the determinant ratio and $\tau$ the corresponding trace defined
in \cite{MagnanoSokolowski}.  The form \eqref{eq:LM} separates Einstein
gravity from a minimally coupled nonlinear massive spin-$2$ field.

For linearly traceless external fields, the determinant and inverse-
determinant contributions to the apparent $\operatorname{tr}\Phi^3$
potential cancel exactly:
\begin{equation}\label{eq:potentialcancel}
V(\Phi)=\frac{m^2}{8}\operatorname{tr}\Phi^2+O(\Phi^4).
\end{equation}
The nonzero $MMM$ amplitude below comes entirely from $K(Q)$.

\subsection{Exact three-point amplitudes}

Take all momenta incoming.  For the $MMM$ vertex use
\begin{equation}\label{eq:mmmkin}
P=(1,0,0,0),\qquad
p=\left(-\frac12,0,0,\frac{i\sqrt3}{2}\right),\qquad
q=\left(-\frac12,0,0,-\frac{i\sqrt3}{2}\right).
\end{equation}
Then $P+p+q=0$ and $P^2=p^2=q^2=-1$.  With normalized TT
polarizations $(M_+,M_+,M_0)$, direct expansion of \eqref{eq:LM} gives
\begin{equation}\label{eq:mmmvalue}
\boxed{\mathcal A_3(M_+,M_+,M_0)=-\frac{\sqrt6}{8}.}
\end{equation}

For $MMh$, take
\begin{equation}\label{eq:mmhkin}
-P=(-1,0,0,0),\qquad
r=(1,-1,-i,0),\qquad
k=(0,1,i,0),
\end{equation}
so $-P+r+k=0$, $r^2=-1$, and $k^2=0$.  For the polarization choice in
Appendix~\ref{app:eframe},
\begin{equation}\label{eq:mmhvalue}
\boxed{\mathcal A_3(M_+,M,h)=-\frac{\sqrt2}{8}.}
\end{equation}
The graviton Ward identity vanishes for an arbitrary symbolic vector
$\xi_\mu$ under
$\varepsilon^h_{\mu\nu}\mapsto k_{(\mu}\xi_{\nu)}$, and exchange of
the two massive legs leaves the coefficient unchanged.

\subsection{Complete massive-pole residue}

Glue \eqref{eq:mmmvalue} and \eqref{eq:mmhvalue} across $P$.  A complete
orthonormal rest-frame basis contains the polarizations
$+,\times,0,xz,yz$.  The exact products are
\begin{center}
\begin{tabular}{c|c|c|c}
\hline
$\lambda$ & $\mathcal A_{MMM}$ & $\mathcal A_{MMh}$ & product\\
\hline
$+$ & $-\sqrt6/8$ & $-\sqrt2/8$ & $\sqrt3/32$\\
$\times$ & $0$ & $0$ & $0$\\
$0$ & $0$ & $\sqrt6/8$ & $0$\\
$xz$ & $0$ & $0$ & $0$\\
$yz$ & $0$ & $i\sqrt2/4$ & $0$\\
\hline
\end{tabular}
\end{center}
Therefore the complete polarization numerator is
\begin{equation}\label{eq:resnum}
N=\sum_{\lambda=1}^{5}
\mathcal A_3(M,M,M_P^\lambda)
\mathcal A_3(M_{-P}^\lambda,M,h)
=\frac{\sqrt3}{32}.
\end{equation}
The TT inverse kernel in the same action convention is
\begin{equation}\label{eq:inversekernel}
K_M(P)=\frac{P^2+M^2}{4}.
\end{equation}
Consequently
\begin{equation}\label{eq:residue}
\boxed{
\Res_{P^2=-M^2}\mathcal A_4(M,M,M,h)=\frac{\sqrt3}{8}\ne0 .}
\end{equation}
An overall rescaling of the action rescales the displayed residue but not
its nonvanishing.  A local quartic contact term is regular at this pole
and cannot cancel it.

\begin{proposition}[Factorization nonvanishing]\label{prop:factorization}
The tree-level four-point amplitude tensor for $MM\to Mh$ is not
identically zero on the relevant complexified on-shell component.
\end{proposition}

\begin{proof}
If the rational amplitude vanished identically, every factorization
residue would vanish.  Equation~\eqref{eq:residue} is the complete
massive-polarization residue and is nonzero.
\end{proof}

\section{The real physical shell}\label{sec:realshell}

\subsection{Open kinematics}

In the center-of-mass frame, with signature $(+,-,-,-)$ restored,
\begin{equation}\label{eq:comkin}
|\mathbf k|=\frac{s-M^2}{2\sqrt s},
\qquad
E_{M,\mathrm{out}}=\frac{s+M^2}{2\sqrt s},
\qquad s\ge4M^2.
\end{equation}
Thus $MM\to Mh$ is physically open on a continuum shell, unlike the
cubic $MMM$ channel.

\begin{lemma}[Complex nonvanishing and a real component]
\label{lem:zariski}
Fix analytic physical polarization sections on an irreducible component
of the complexified four-point mass shell.  The nonsingular real
$MM\to Mh$ region contains a Euclidean-open center-of-mass chart and is
Zariski dense in the complex component containing it.  Hence a rational
tree amplitude that vanishes on an open real subset vanishes identically
on that component and has zero factorization residues.
\end{lemma}

\begin{proof}
Away from exceptional polarization charts, the real physical region is
parametrized by $s>4M^2$, $-1<\cos\theta<1$, and an azimuth.  Its real
points contain an open set of maximal real dimension in the relevant
algebraic component.  If $\mathcal A=P/Q$ is regular there and vanishes
on a real open set, real analyticity makes $P$ vanish there.  The
identity theorem, equivalently Zariski density of that real locus, makes
$P$ vanish on the component.  Then all residues vanish, contrary to
\eqref{eq:residue}.
\end{proof}

The next subsection supplies a stronger direct bridge: an exact nonzero
point already on the real shell.  It makes Lemma~\ref{lem:zariski}
logically redundant for the obstruction theorem, while retaining it as
an independent explanation of why the complex factorization rail is
physically relevant.

\subsection{Interior rational certificate}

Set $M=1$.  With $p_1,p_2$ incoming and $p_3,k$ outgoing, take
\begin{align}\label{eq:rationalpoint}
p_1&=\left(\frac54,0,0,\frac34\right),
&p_2&=\left(\frac54,0,0,-\frac34\right),\nonumber\\
p_3&=\left(\frac{29}{20},-\frac{21}{25},0,-\frac{63}{100}\right),
&k&=\left(\frac{21}{20},\frac{21}{25},0,\frac{63}{100}\right).
\end{align}
Then
\begin{equation}
p_1+p_2=p_3+k,\qquad p_i^2=1,\qquad k^2=0,
\end{equation}
and
\begin{equation}
s=\frac{25}{4}M^2,
\qquad \cos\theta=\frac35.
\end{equation}
This point is interior and noncollinear.  The planar kinematics has a
$y$-reflection selection rule; the verification therefore uses explicit
parity-even real linear polarizations rather than interpreting symmetry
zeros as dynamical cancellations.

\begin{cproposition}[Real-shell certificate, G15]\label{cprop:g15}
For the real polarization tensors listed in Appendix~\ref{app:fourpoint},
the complete reduced tree amplitude at \eqref{eq:rationalpoint} is
\begin{equation}\label{eq:AK}
\boxed{\AK(MM\to Mh)=\frac{7881241032}{5584765625}\ne0.}
\end{equation}
It is real, gauge independent, and nonsingular at this point.  Therefore
it is nonzero on a nonempty open neighborhood of the real physical shell.
\end{cproposition}

\section{Full fourth-order calculation and the metric obstruction}
\label{sec:full}

The calculation of Proposition~\ref{cprop:g15} is performed directly in
the original fourth-order variables of \eqref{eq:action}, independently
of the Einstein-frame factorization rail.

\subsection{Perturbiner and bordered propagator}

For four plane waves,
\begin{equation}
g_{\mu\nu}=\eta_{\mu\nu}
+\sum_{a=1}^{4}\lambda_a\varepsilon^{(a)}_{\mu\nu}e^{ip_a\cdot x},
\end{equation}
the wave algebra retains only multilinear subsets of the labels and
implements $\partial_\mu\mapsto iP_\mu$ automatically.  The coefficient
of $\lambda_1\lambda_2\lambda_3\lambda_4$ in the action gives the exact
quartic contact.  The same engine produces the quadratic $10\times10$
operator $K(P)$ and cubic current $J_A$ in the symmetric-tensor basis.

No symbolic closed propagator is assumed.  In each $s,t,u$ channel one
solves the exact de Donder-bordered system
\begin{equation}\label{eq:bordered}
\begin{pmatrix}K(P)&C(P)^T\\ C(P)&0\end{pmatrix}
\begin{pmatrix}X\\ \lambda\end{pmatrix}
=\begin{pmatrix}-J^{(cd)}\\0\end{pmatrix},
\qquad
\mathcal A_{\rm exch}=J^{(ab)T}X,
\end{equation}
where
\begin{equation}
C_\mu(X)=P^\alpha X_{\alpha\mu}-\frac12P_\mu X^\alpha{}_\alpha.
\end{equation}
The exact field-equation residual and $CX$ vanish in all three channels.
Replacing $C$ by the axial constraint $n^\mu X_{\mu\nu}=0$ leaves the
total amplitude unchanged.

\subsection{Contact and exchange decomposition}

At the rational point and polarization choice, the four contributions
are
\begin{align}\label{eq:pieces}
\mathcal A_{\rm contact}&=
\frac{8364585636057}{410644531250},\nonumber\\
\mathcal A_s&=-\frac{836906329}{45312500},\nonumber\\
\mathcal A_t&=-\frac{5142800561}{111695312500},\nonumber\\
\mathcal A_u&=-\frac{90855828716}{205322265625}.
\end{align}
Their exact sum is \eqref{eq:AK}.  The contact term alone is not gauge
invariant.  Replacing the external graviton by
$k_{(\mu}\xi_{\nu)}$ makes the \emph{complete} contact-plus-exchange
sum vanish for an arbitrary exact rational test vector; shifting a
physical representative by the same gauge tensor leaves the amplitude
unchanged.  Initial-massive Bose exchange is exact.  The two-point
coefficient against every symmetric-tensor basis direction vanishes for
each external on-shell polarization, excluding equation-of-motion
contamination.

Reversing all external momenta and conjugating the real polarization
sections reproduces the contact and each exchange separately.  The
quadratic kernels are recomputed and satisfy $K(-P)=K(P)$ exactly.  Thus
the reverse process entering the physical adjoint is evaluated, not
inferred.

\subsection{Quarter-turn of the complete operator}

For the contact, three external massive fields give $(-i)^3=i$.  An
exchange through an internal massless line has endpoint phase
$(-i)^2(-i)=i$.  An exchange through an internal massive line has
endpoint phase $(-i)^3(-i)^2=-i$, while the quarter-turned massive
inverse kernel contributes $-1$, again giving $i$.  Consequently every
piece of \eqref{eq:pieces} obeys Theorem~\ref{thm:phase}, and
\begin{equation}\label{eq:tplus}
\langle Mh|T_{2,+}|MM\rangle=i\AK,
\qquad
\langle MM|T_{2,+}|Mh\rangle=i\AK.
\end{equation}

Both states have exact free energy
\begin{equation}
E_{MM}=\frac54+\frac54=\frac52
=\frac{29}{20}+\frac{21}{20}=E_{Mh}.
\end{equation}
On the ordered shell basis $(|MM\rangle,|Mh\rangle)$,
\begin{equation}\label{eq:Tblock}
T_{2,+}=i\AK\,\sigma_x,
\end{equation}
and hence
\begin{equation}\label{eq:obstruction}
\boxed{
\PiE\bigl(T_{2,+}^{\dg}-T_{2,+}\bigr)
=-2i\AK\,\sigma_x
=-\frac{15762482064\,i}{5584765625}\,\sigma_x\ne0.}
\end{equation}

\subsection{First-order metric ambiguity}

The solution $R_1$ of the first equation in \eqref{eq:defcomplex} is
defined up to a Hermitian $G$ with $[G,h_0]=0$.  Such a change shifts the
second-order source by $\tfrac12[G,v_1+v_1^{\dg}]$.  On the nonsingular
real physical shell the complete connected cubic block is zero: $Mhh$
vanishes by Lemma~\ref{lem:onemassive}; $MMM$ and $MMh$ have no nonsoft
real $1\leftrightarrow2$ shell; and real $hhh$ three-point kinematics is
collinear with zero physical helicity amplitude.  Therefore
\begin{equation}\label{eq:ambiguity}
P_E(v_1+v_1^{\dg})P_E=0,
\qquad
P_E[G,v_1+v_1^{\dg}]P_E=0.
\end{equation}
The obstruction cannot be moved or canceled by choosing another analytic
first-order metric representative.

\begin{theorem}[Second-order positive-metric obstruction]
\label{thm:positiveobstruction}
The Lorentz-covariant positive free completion of scalar-free
Einstein--Weyl gravity admits no analytic second-order metric deformation
compatible with the full interaction on the physical $MM\to Mh$ shell.
The obstruction is nonzero at \eqref{eq:rationalpoint} and on a nonempty
open neighborhood of that point.
\end{theorem}

\begin{proof}
Equation~\eqref{eq:obstruction} lies in $\ker\ad_{h_0}$ and is nonzero.
No $R_2$ can solve $[h_0,R_2]$ equal to this source.  Equation
\eqref{eq:ambiguity} makes the class independent of first-order
commutant freedom.  The exact Ward, external-EOM, reverse-process, and
internal-gauge checks show that the displayed matrix element is a reduced
physical one.  Since it is a nonsingular real-analytic amplitude and is
nonzero at an interior point, it remains nonzero on an open neighborhood.
\end{proof}
\section{Conventional covariant Krein grading}\label{sec:grading}

The positive obstruction does not by itself decide whether an indefinite
completion can still be formulated.  The narrower question addressed
here is whether the conventional gravitational Krein grading can render
the offending block null, as a continuous one-sided charge does at the
perfect-square scalar boundary.

\subsection{One-particle commutant}

Let $J^{(1)}$ be a nondegenerate Hermitian involution on the physical
one-particle cohomology \eqref{eq:oneparticle}.  We require it to commute
with the proper-orthochronous Poincar\'e representation and to agree with
the free Krein real form of Paper~IV.

The mass Casimir separates the massless and massive orbits, so no
covariant endomorphism mixes $\mathcal H_h$ and $\mathcal H_M$.  On the
massive fiber, the exact spin-$2$ $SO(3)$ commutant is one-dimensional.
There is one precision on the massless fiber: under the
proper-orthochronous group, helicity $+2$ and helicity $-2$ are
inequivalent complex representations, so covariance alone permits
separate scalars.  The complete commutant is therefore
\begin{equation}\label{eq:commutant}
J^{(1)}=\epsilon_+I_{+2}\oplus\epsilon_-I_{-2}
\oplus\epsilon_MI_5,
\qquad \epsilon_\pm,\epsilon_M\in\{\pm1\}.
\end{equation}
Parity, or compatibility with the real gravitational field, exchanges
the two massless helicities and imposes $\epsilon_+=\epsilon_-$.  Even
without adding parity to the covariance group, agreement with the free
signature \eqref{eq:freesignature} fixes
\begin{equation}\label{eq:epsfix}
\epsilon_+=\epsilon_-=+1,
\qquad \epsilon_M=-1.
\end{equation}

\begin{proposition}[One-particle grading classification]
\label{prop:oneparticlegrading}
Within the stated covariant class, the free real form uniquely fixes
\begin{equation}
J^{(1)}=I_2\oplus(-I_5).
\end{equation}
No linear covariant grading can assign different signs within the
massive spin-$2$ multiplet or mix the massive and massless shells.
\end{proposition}

\subsection{Cluster-factorizing Fock lift}

Let $j(n_h,n_M)\in\{\pm1\}$ be the sign assigned to an asymptotic sector
with $n_h$ massless and $n_M$ massive particles.  Tensor multiplicativity
and cluster decomposition require
\begin{equation}\label{eq:monoid}
j(n_h+n'_h,n_M+n'_M)
=j(n_h,n_M)j(n'_h,n'_M),
\qquad j(0,0)=1.
\end{equation}
A homomorphism from the free commutative monoid $\mathbb N^2$ is fixed by
its values on the two generators.  Equations \eqref{eq:epsfix} and
\eqref{eq:monoid} therefore give
\begin{equation}\label{eq:fockgrading}
\boxed{J_F=(-1)^{N_M}.}
\end{equation}
Assigning independent signs to unrelated multiparticle sectors would
violate \eqref{eq:monoid} and hence cluster factorization.

At the G17 shell,
\begin{equation}\label{eq:statesigns}
J_F|MM\rangle=+|MM\rangle,
\qquad
J_F|Mh\rangle=-|Mh\rangle.
\end{equation}
The transition crosses the two Krein signatures.

\section{Krein visibility of the obstruction}\label{sec:visibility}

For an operator on the positive-frame Hilbert representatives, define
the Krein adjoint
\begin{equation}\label{eq:sharp}
X^\sharp=J_FX^{\dg}J_F.
\end{equation}
On the minimal ordered basis $(|MM\rangle,|Mh\rangle)$,
$J_{\min}=\operatorname{diag}(1,-1)$.

Consider first the forward block
\begin{equation}\label{eq:Xforward}
X=i\AK\,|Mh\rangle\langle MM|.
\end{equation}
It is $J_F$-odd:
\begin{equation}
J_FXJ_F=-X.
\end{equation}
Its Krein adjoint and quadratic product are
\begin{align}\label{eq:Xsharp}
X^\sharp&=-i\AK\,J_F|MM\rangle\langle Mh|J_F
=i\AK\,|MM\rangle\langle Mh|,\nonumber\\
X^\sharp X&=-\AK^2|MM\rangle\langle MM|.
\end{align}
Consequently
\begin{equation}\label{eq:Xtrace}
\boxed{
\Tr(X^\sharp X)=-\AK^2
=-\frac{62113960204480425024}{31189607086181640625}\ne0.}
\end{equation}
The sign is indefinite, but the operator is not null.

For the complete positive-metric obstruction
\begin{equation}
\mathcal O=-2i\AK\sigma_x,
\end{equation}
one obtains
\begin{equation}\label{eq:Osharp}
\boxed{
\mathcal O^\sharp=\mathcal O,
\qquad
\mathcal O^\sharp\mathcal O=-4\AK^2I_2,
\qquad
\Tr(\mathcal O^\sharp\mathcal O)=-8\AK^2\ne0.}
\end{equation}
In the exact rational convention,
\begin{equation}
-8\AK^2
=-\frac{496911681635843400192}{31189607086181640625}.
\end{equation}

\begin{proposition}[Conventional Krein non-nullity]
\label{prop:nonnull}
The physical $MM\to Mh$ block and the complete G17 obstruction are
non-null for the unique conventional covariant, cluster-factorizing
fundamental symmetry \eqref{eq:fockgrading}.  The obstruction is
Krein-self-adjoint, indefinite, and visible to the finite-block quadratic
trace.
\end{proposition}

The algebraic distinction from a one-sided continuous grading is simple.
For $\mathbb Z_2$,
\begin{equation}
1+1=0\pmod2.
\end{equation}
The adjoint of an odd operator is odd and their product is even.  Thus
oddness supplies no charge-null lemma:
\begin{equation}\label{eq:oddnotnull}
\boxed{J_F\text{-odd}\ \not\Longrightarrow\ \text{null}.}
\end{equation}
By contrast, at the scalar $O(1,1)$ boundary a strictly one-sided
nonzero continuous charge remains nonzero after the relevant product,
and an invariant graded trace can kill it \cite{paperV,BT2026}.

\begin{remark}[What is not ruled out]
Equation~\eqref{eq:Osharp} does not say that a Krein-self-adjoint
interacting Hamiltonian cannot exist.  Indeed, it displays
Krein-self-adjointness of the minimal obstruction block.  What fails is
the stronger proposal that the conventional massive-number grading makes
the branch-changing component null or supplies a positive generalized
Born contribution.  Indefinite pseudo-unitary evolution and a positive
probability calculus are different questions.
\end{remark}

\section{No scalar-like charge-null mechanism}\label{sec:charge}

Could a continuous internal charge make the obstruction one-sided?
Poincar\'e covariance first restricts any linear
charge on the irreducible massive fiber to one scalar $q_M$.  The free
real graviton is neutral, $q_h=0$.  The verified nonzero interaction
vertices impose
\begin{equation}\label{eq:charges}
3q_M=0\quad(MMM),
\qquad
2q_M+q_h=2q_M=0\quad(MMh).
\end{equation}

\begin{proposition}[Uniform abelian charge no-go]\label{prop:chargenogo}
There is no nontrivial uniform abelian charge assignment to the physical
massive spin-$2$ irrep, with neutral graviton, that is respected by the
Einstein--Weyl interaction.
\end{proposition}

\begin{proof}
Equations \eqref{eq:charges} hold in any abelian charge group.  Since
$q_M=3q_M-2q_M$, they imply $q_M=0$.  The conclusion does not rely on
torsion-freeness; $\gcd(2,3)=1$ eliminates finite torsion charges as well.
\end{proof}

The $MMM$ vertex has no real cubic decay shell but is a nonzero algebraic
interaction vertex, so it already breaks massive-number parity at the
action level.  More decisively, the real G17 process itself changes
$N_M$ from $2$ to $1$ and proves
\begin{equation}\label{eq:paritybroken}
[(-1)^{N_M},S]\ne0
\end{equation}
on the physical scattering shell.  Therefore the ingredients used to
prove factorization simultaneously exclude an $O(1,1)$-like charge
replacement.

\section{BRST cohomology}\label{sec:brst}

The higher-derivative massive branch is a physical reduced mode and must
not be confused with Faddeev--Popov ghosts.  The external states in G17
are TT massive polarizations and TT massless graviton polarizations.
The exact Ward identity, invariance under changes of gauge representative,
axial/de Donder equality, and external inverse-propagator
zeros show that the amplitude descends to reduced physical BRST
cohomology.

Let $Q$ denote the nilpotent BRST charge and let $|i\rangle,|f\rangle$
be closed physical representatives.  If an operator were BRST-exact,
\begin{equation}
\mathcal O=\{Q,Y\},
\end{equation}
then, under the standard decoupling assumptions,
\begin{equation}\label{eq:brstzero}
\langle f|\mathcal O|i\rangle
=\langle f|QY|i\rangle+\langle f|YQ|i\rangle=0.
\end{equation}
The matrix element \eqref{eq:obstruction} is nonzero.

\begin{proposition}[BRST survival]\label{prop:brst}
The G17 obstruction defines a nonzero operator on physical BRST
cohomology and is not BRST-exact:
\begin{equation}
[\mathcal O]\ne0.
\end{equation}
It cannot be placed in the standard BRST-null ideal.
\end{proposition}

Appendix~\ref{app:brst} records the finite-complex algebra used in G18.
In a splitting into physical cohomology and a contractible doublet, the
physical-to-physical block of $\{Q,Y\}$ vanishes for a completely general
$Y$, while the G17 block survives unchanged.

\section{Master theorem}\label{sec:master}

\begin{theorem}[Interacting completion obstruction]\label{thm:master}
Let scalar-free Einstein--Weyl gravity \eqref{eq:action} be in its regular
split phase about flat space.  Use the covariant free completion
classes of Paper~IV and the analytic metric-deformation class of
Paper~V.  Then:
\begin{enumerate}
\item the Lorentz-covariant positive pseudo-Hermitian completion is
protected at first perturbative order, $\ob_+^{(1)}=0$;
\item at second order the physical process $MM\to Mh$ generates the
nonzero shell obstruction \eqref{eq:obstruction};
\item hence the canonical covariant positive metric has no analytic
second-order deformation compatible with the full interaction;
\item a conventional Poincar\'e-covariant, real-field-compatible,
BRST-compatible, nondegenerate and cluster-factorizing Krein grading
agreeing with the free real form has the unique Fock lift
$J_F=(-1)^{N_M}$;
\item for this grading the obstruction is non-null,
$\Tr(\mathcal O^\sharp\mathcal O)=-8\AK^2\ne0$, and survives physical
BRST cohomology; and
\item the nonzero $MMM$ and $MMh$ vertices prohibit a nontrivial uniform
abelian charge grading analogous to the scalar $O(1,1)$ mechanism.
\end{enumerate}
Therefore neither an analytic deformation of the canonical positive
metric nor the conventional massive-number null/parity mechanism
supplies an interacting positive-probability completion of the full
theory in the stated class.
\end{theorem}

\begin{proof}
Item 1 is Theorem~\ref{thm:cubicprotection}.  Items 2 and 3 are
Theorem~\ref{thm:positiveobstruction}, independently supported by the
factorization Proposition~\ref{prop:factorization}.  Item 4 is
Proposition~\ref{prop:oneparticlegrading} together with the monoid
argument \eqref{eq:monoid}.  Item 5 is
Propositions~\ref{prop:nonnull} and \ref{prop:brst}.  Item 6 is
Proposition~\ref{prop:chargenogo}.
\end{proof}

\section{Scope and nonclaims}\label{sec:scope}

The hypotheses in Theorem~\ref{thm:master} are part of the result.  The
paper rules out:
\begin{itemize}
\item analytic deformation of the canonical Lorentz-covariant positive
free metric;
\item a conserved naive massive-number ghost parity;
\item conventional covariant linear one-particle fundamental symmetries
that agree with the free real form but differ on selected massive
helicities;
\item non-multiplicative independent Fock-sector signs, once cluster
factorization is required;
\item a uniform abelian continuous- or torsion-charge null mechanism
compatible with the verified vertices; and
\item standard BRST-exact nullity of the physical obstruction.
\end{itemize}

It does \emph{not} prove:
\begin{itemize}
\item that no nonperturbative or differently pointed positive completion
exists;
\item that the full interacting Hamiltonian has complex energies;
\item that every nonlocal or non-factorizing grading fails;
\item that no singular $M\to0$ or conformal-boundary completion can be
defined;
\item that a Maldacena-type boundary truncation of the solution algebra is
inconsistent \cite{Maldacena}; or
\item that the result extends unchanged to quadratic gravities with a
scalaron, matter, cosmological constant, or a different background.
\end{itemize}

Any surviving null-sector construction must therefore leave the
conventional class.  Possibilities include a grading singular at the
conformal boundary, a nonlinear mixing of asymptotic particle sectors, a
nonlocal observable algebra, a failure of ordinary cluster
factorization, or boundary conditions that remove part of the
fourth-order solution space.  These are qualitatively different from a
straightforward gravitational lift of the Bateman--Turok mechanism.

\section{Discussion}\label{sec:discussion}

\subsection{The same order pattern, for different reasons}

Paper~V and the present gravity calculation share the sequence
\begin{equation}
\text{first-order protection}\longrightarrow
\text{second-order scattering obstruction}.
\end{equation}
The agreement in order should not obscure the difference in structure.
The perfect-square scalar is protected by its special cubic kinematic
factor and reality assignments.  Gravity is protected because the
Einstein solution space is an exact nonlinear subsector and the remaining
odd cubic vertex is kinematically closed.  In the scalar, an enhanced
$O(1,1)$ charge appears exactly at the perfect-square boundary and can
sort an obstruction into a one-sided null ideal.  In gravity, the $MMM$
and $MMh$ interactions remove every uniform abelian charge before the
four-point obstruction is considered.

\subsection{Real forms become dynamical}

At the free level a real form selects which complex solutions are called
real and which adjoint defines observables.  The choice can appear
kinematic.  The phase theorem makes its interacting content explicit:
every physical process with odd external massive number crosses the two
reality assignments and becomes anti-Hermitian in the positive frame.
The decisive diagnostic is therefore
\begin{equation}
\boxed{\text{Which physical on-shell processes change the real-form
sector?}}
\end{equation}
For Einstein--Weyl gravity the first answer is $MM\to Mh$.

\subsection{Why the factorization and physical rails are both useful}

The Einstein-frame residue \eqref{eq:residue} is compact, independent,
and immune to cancellation by local contact terms.  It proves a genuine
interaction before the large quartic calculation is attempted.  The
original-variable physical point then supplies what factorization alone
does not: an exact real-shell value, the full gauge cancellation, the
physical adjoint, and the deformation-complex matrix element.  Agreement
between the two rails is substantially stronger than either calculation
in isolation.

\subsection{What Krein visibility means}

The negative traces in \eqref{eq:Xtrace} and \eqref{eq:Osharp} are not
advertised as negative probabilities.  They show that the conventional
quadratic Krein form sees the block: it is neither zero nor killed by a
grading selection rule.  A viable interacting probability construction
would need extra structure beyond the free fundamental symmetry.  The
scalar boundary supplies such extra structure through a continuous
one-sided charge; the verified Einstein--Weyl vertices do not.

The immediate research problem is consequently no longer another scalar
doubling.  It is to understand whether conformal-boundary conditions,
changed observable algebras, or genuinely nonlocal gravitational
structures can evade the hypotheses of Theorem~\ref{thm:master} without
destroying covariance and factorization.

\begin{thebibliography}{16}

\bibitem{Stelle}
K.~S.~Stelle,
\emph{Renormalization of higher-derivative quantum gravity},
Phys.\ Rev.\ D \textbf{16} (1977) 953.

\bibitem{MagnanoSokolowski}
G.~Magnano and L.~M.~Soko\l{}owski,
\emph{Nonlinear massive spin-two field generated by higher derivative
gravity},
Ann.\ Phys.\ \textbf{306} (2003) 1--35,
arXiv:gr-qc/0209022.

\bibitem{LiuLuPope}
H.-S.~Liu, H.~L\"u, C.~N.~Pope and J.~F.~V\'azquez-Poritz,
\emph{Not conformally-Einstein metrics in conformal gravity},
arXiv:1303.5781.

\bibitem{KOS}
S.~Kamefuchi, L.~O'Raifeartaigh and A.~Salam,
\emph{Change of variables and equivalence theorems in quantum field
theories},
Nucl.\ Phys.\ \textbf{28} (1961) 529--549.

\bibitem{BM2008}
C.~M.~Bender and P.~D.~Mannheim,
\emph{No-ghost theorem for the fourth-order derivative Pais--Uhlenbeck
oscillator model},
Phys.\ Rev.\ Lett.\ \textbf{100} (2008) 110402.

\bibitem{BM2008PRD}
C.~M.~Bender and P.~D.~Mannheim,
\emph{Exactly solvable PT-symmetric Hamiltonian having no Hermitian
counterpart},
Phys.\ Rev.\ D \textbf{78} (2008) 025022, arXiv:0804.4190.

\bibitem{BT2026}
S.~Bateman and N.~Turok,
\emph{Escape from Ostrogradsky via hidden ghost parity},
arXiv:2607.00096.

\bibitem{Maldacena}
J.~Maldacena,
\emph{Einstein gravity from conformal gravity},
arXiv:1105.5632.

\bibitem{AzizovIokhvidov}
T.~Ya.~Azizov and I.~S.~Iokhvidov,
\emph{Linear Operators in Spaces with an Indefinite Metric},
Wiley, 1989.

\bibitem{paperI}
\emph{Canonical positive symplectic diagonalization of the
Pais--Uhlenbeck oscillator}, companion paper (Paper~I), 2026.

\bibitem{paperII}
\emph{The Pais--Uhlenbeck metric as a minimum-distortion principle, and
the representation problem for the fourth-order field}, companion paper
(Paper~II), 2026.

\bibitem{paperIII}
\emph{The universal vacuum of the fourth-order scalar field: metric
orbits, Fock sectors, and the Krein boundary}, companion paper
(Paper~III), 2026.

\bibitem{paperIV}
\emph{Gauge reduction and the completion problem in fourth-order
gravity: PU pairing, covariant real forms, and the conformal Jordan
boundary}, companion paper (Paper~IV), 2026.

\bibitem{paperV}
\emph{Interaction obstructions, resonant PT breaking, and doubled Jordan
symmetry in fourth-order theories}, companion paper (Paper~V), 2026.

\end{thebibliography}

\appendix

\section{Einstein consistent truncation and LSZ}
\label{app:truncation}

This appendix spells out the hypotheses behind
Lemma~\ref{lem:onemassive}.  For the scalar-free combination in
\eqref{eq:action}, the field equation is the Einstein tensor plus a
multiple of the Bach tensor.  In four dimensions the Bach tensor may be
written
\begin{equation}
B_{\mu\nu}=\left(\nabla^\rho\nabla^\sigma
+\frac12R^{\rho\sigma}\right)C_{\mu\rho\nu\sigma}
\end{equation}
up to a convention-dependent overall sign.  It vanishes on every
Einstein metric \cite{LiuLuPope}.  In the flat-background problem the
relevant exact subsector is Ricci-flat vacuum Einstein gravity.

On the regular branch of the Legendre transformation from Ricci-tensor
gravity to its second-order form, introduce the Einstein metric and a
symmetric massive field $\Phi_{\mu\nu}$.  The transformation is algebraic
in the auxiliary variables and locally invertible wherever its determinant
does not vanish \cite{MagnanoSokolowski}.  Linearization about the common
flat vacuum diagonalizes the asymptotic equations into
\begin{equation}
\mathcal K_hh=0,
\qquad
\mathcal K_MM=0,
\end{equation}
with $p^2=0$ and $p^2=M^2$, respectively.  Thus $h$ and $M$ used in the
main text are LSZ mass eigenfields, not arbitrary components of the
original metric perturbation.

Let $E_M[h,M]=0$ be the exact massive-field equation in this frame.  The
consistent truncation states
\begin{equation}\label{eq:truncapp}
E_M[h_E,0]=0
\end{equation}
for every nonlinear Einstein solution $h_E$.  Expanding about flat space,
\begin{equation}
E_M[h,M]=\mathcal K_MM+
\sum_{n\ge2}J_M^{(n)}[h^n]
+\sum_{r\ge1,s\ge1}J_{r,s}[M^rh^s].
\end{equation}
Equation \eqref{eq:truncapp} does not require each off-shell functional
$J_M^{(n)}$ to vanish.  It requires its value on a full perturbative
Einstein solution to cancel, including nonlinear corrections to that
solution.  LSZ amputation removes precisely the external inverse-
propagator and field-redefinition terms, leaving
\begin{equation}
\mathcal A(M,h_1,\ldots,h_n)=0
\end{equation}
for physical on-shell Einstein gravitons.  This is the standard
consistent-truncation selection rule.

The perturbiner supplies an explicit cubic check.  In signature
$(+,-,-,-)$, take all momenta incoming,
\begin{equation}
p_M=(1,0,0,0),\qquad
p_2=\left(-\frac12,0,0,-\frac12\right),\qquad
p_3=\left(-\frac12,0,0,\frac12\right).
\end{equation}
They obey $p_M^2=1$, $p_2^2=p_3^2=0$ and sum to zero.  The exact action
coefficient vanishes for all five massive TT polarizations and the four
products of the two graviton helicities.  Replacing either graviton
polarization by $p_{(\mu}\xi_{\nu)}$ also gives zero.

The selection rule is tree-level and concerns external mass eigenfields.
It does not state that the original fourth-order metric action lacks
terms apparently linear in a chosen massive component, nor does it
exclude loop effects, anomalies, or a singular failure of the
mass-eigenfield map at the conformal boundary.

\section{Einstein-frame cubic expansion}
\label{app:eframe}

We record enough of the independent G14 calculation to make the exact
numbers reproducible.  This appendix uses signature $(-,+,+,+)$ and
$M=1$.  It matches
\path{symbolic/verify_gravity_factorization.py}.

For a symmetric matrix $\Phi$, write
\begin{equation}
X=g^{-1}\Phi,\qquad
t_n=\operatorname{tr}(X^n),\qquad
A=\sqrt{\det(I+X)}.
\end{equation}
Through cubic order,
\begin{align}
A={}&1+\frac12t_1+\frac18t_1^2-\frac14t_2
+\frac1{48}t_1^3-\frac18t_1t_2+\frac16t_3+O(\Phi^4),\nonumber\\
A^{-1}={}&1-(A-1)+(A-1)^2-(A-1)^3+O(\Phi^4).
\end{align}
In the potential of \eqref{eq:LM}, $\tau=4+t_1$ and the relevant
quadratic contraction is $4+2t_1+t_2$.  Substitution gives, when
$t_1=0$,
\begin{equation}
A^{-1}\left(-\tau^2+\psi_{\mu\nu}\psi^{\mu\nu}
+6A\tau-12A^2\right)=t_2+O(\Phi^4),
\end{equation}
with zero $t_3$ coefficient.  This is the potential cancellation
\eqref{eq:potentialcancel}.

At first order in one plane wave $(p,E)$,
\begin{equation}
Q^{(1)\alpha}{}_{\mu\nu}(E,p)
=\frac{i}{2}\eta^{\alpha\beta}
\left(p_\mu E_{\beta\nu}+p_\nu E_{\beta\mu}
-p_\beta E_{\mu\nu}\right).
\end{equation}
The quadratic correction $Q^{(2)}$ comes from the inverse-$\psi$
factor.  The cubic $MMM$ coefficient is the sum of all contractions
$K(Q^{(1)},Q^{(2)})+K(Q^{(2)},Q^{(1)})$, with the factor $1/2$ in
\eqref{eq:LM}.  Exact substitution of \eqref{eq:mmmkin} and normalized
TT tensors gives \eqref{eq:mmmvalue}.

For $MMh$, the script expands $g^{-1}$, $\psi^{-1}$, the connection of
$g$, and $Q$ through the orders required for the coefficient of three
independent wave labels.  At \eqref{eq:mmhkin}, take
\begin{equation}
v_r=(1,-1,0,0),\qquad e_h=(1,0,0,1),
\qquad E_r=v_r^\flat\otimes v_r^\flat,
\qquad H=e_h^\flat\otimes e_h^\flat.
\end{equation}
Together with the internal $+$ polarization at rest these tensors are
transverse and traceless, and the coefficient is \eqref{eq:mmhvalue}.
For
\begin{equation}
H^{\rm gauge}_{\mu\nu}=k_\mu\xi_\nu+k_\nu\xi_\mu
\end{equation}
with four independent symbolic components of $\xi$, the coefficient
vanishes identically.

Finally, the massive TT two-point coefficient for waves $(p,E)$ and
$(-p,E)$ with $E\cdot E=1$ is
\begin{equation}
L_M^{(2)}=\frac{p^2+1}{4}.
\end{equation}
This fixes the factor of four between the polarization numerator
$\sqrt3/32$ and pole residue $\sqrt3/8$; it is not inferred from a
unit-residue completeness convention.

\section{Polarization basis and factorization sum}
\label{app:polarizations}

At the rest momentum $P=(1,0,0,0)$ in the Einstein-frame signature,
let $e_x,e_y,e_z$ be unit spatial vectors.  A normalized massive basis is
\begin{align}
E_+&=\frac1{\sqrt2}(e_xe_x-e_ye_y),
&E_\times&=\frac1{\sqrt2}(e_xe_y+e_ye_x),\nonumber\\
E_0&=\frac1{\sqrt6}(e_xe_x+e_ye_y-2e_ze_z),
&E_{xz}&=\frac1{\sqrt2}(e_xe_z+e_ze_x),\nonumber\\
&&E_{yz}&=\frac1{\sqrt2}(e_ye_z+e_ze_y).
\end{align}
Each is transverse, traceless, and orthonormal under the induced tensor
inner product.  For a massive momentum $(E,0,0,p_z)$, replace $e_z$ by
the normalized longitudinal transverse vector; the same formulas give
the transported helicity basis.

The factorization numerator is basis-independent.  In a nonorthonormal
basis $T_a$ with Gram matrix $G_{ab}=T_a\cdot T_b$, it is
\begin{equation}
N=V_a(G^{-1})^{ab}U_b,
\end{equation}
where $V_a=\mathcal A_3(M,M,T_a)$ and
$U_b=\mathcal A_3(T_b,M,h)$.  The orthonormal basis above reduces this
to the sum displayed after \eqref{eq:resnum}.  All five intermediate
states are included; the result is not one selected nonzero channel.

The original fourth-order perturbiner supplies a separate complex
factorization check at a rational-$i$ point.  It constructs an exact
five-dimensional TT basis at the internal massive momentum, contracts
the two cubic vectors with its Gram inverse, and obtains a nonzero
reduced numerator.  The Einstein-frame calculation then fixes a compact
normalized representative and the explicit inverse-kernel slope.

\section{Physical-shell geometry and the rational point}
\label{app:shell}

For incoming equal masses in the center-of-mass frame,
\begin{equation}
p_1=(E,0,0,p),\qquad p_2=(E,0,0,-p),
\qquad E=\frac{\sqrt s}{2},\quad
p=\frac12\sqrt{s-4M^2}.
\end{equation}
Choose the outgoing graviton at scattering angle $\theta$,
\begin{equation}
k=(\kappa,\kappa\sin\theta,0,\kappa\cos\theta),
\qquad
\kappa=\frac{s-M^2}{2\sqrt s},
\end{equation}
and $p_3=p_1+p_2-k$.  Then $p_3^2=M^2$ and
$p_3^0=(s+M^2)/(2\sqrt s)$.  The domain
\begin{equation}
s>4M^2,\qquad -1<\cos\theta<1
\end{equation}
is a real open chart away from forward/backward and threshold
degeneracies.

At $s=25M^2/4$ one has $E=5M/4$ and
$\kappa=21M/20$.  Taking $\cos\theta=3/5$ and
$\sin\theta=4/5$ gives exactly \eqref{eq:rationalpoint}.  All energy,
momentum, and mass-shell identities are rational.  Because all spatial
$y$ components vanish, reflection $y\mapsto-y$ is a symmetry.  An odd
number of $y$-odd polarization tensors forces an amplitude zero.  The
certificate uses the first nonzero all-even combination.

The exact threshold regression uses two incoming massive particles at
rest and an outgoing graviton of energy $3M/4$.  In the separate
polarization convention of that test,
\begin{equation}
\mathcal A_{\rm threshold}=-\frac{509784}{390625}.
\end{equation}
Threshold is not used for the open-set claim; it guards the implementation
against later changes.

\section{Full four-point engine and G17 expressions}
\label{app:fourpoint}

\subsection{Wave algebra}

A wave-algebra element is a dictionary from subsets of external labels
to sparse tensor components.  Multiplication discards products in which
a wave label repeats, so multilinear truncation is automatic.  A
derivative acting on subset $S$ multiplies by
\begin{equation}
iP_{S\mu}=i\sum_{a\in S}p_{a\mu}.
\end{equation}
From the metric, inverse metric, determinant, Christoffel symbols, and
Ricci tensor, the engine extracts
\begin{equation}
\sqrt{-g}R,\qquad
\sqrt{-g}R_{\mu\nu}R^{\mu\nu},
\qquad
\sqrt{-g}R^2
\end{equation}
at the desired wave subset.  Before any cubic or quartic claim it checks
that the two-point coefficient vanishes on $p^2=0$ and $p^2=M^2$, is
nonzero off shell, and obeys cubic Ward identities.

\subsection{Exact polarization tensors}

For a massive momentum $p$ in the $(t,x,z)$ plane, choose two independent
transverse in-plane vectors $v_1,v_2$.  With
\begin{equation}
\Pi_{\mu\nu}=\eta_{\mu\nu}-\frac{p_\mu p_\nu}{M^2},
\end{equation}
the parity-even TT tensor used in the certificate is
\begin{equation}
E_{12,\mu\nu}=v_{1\mu}v_{2\nu}+v_{2\mu}v_{1\nu}
-\frac23(v_1\cdot v_2)\Pi_{\mu\nu}.
\end{equation}
The three massive external legs use the corresponding $E_{12}$ at
$p_1,p_2,-p_3$.  For the graviton momentum direction proportional to
$(5,4,0,3)$, use transverse vectors
\begin{equation}
e_1=(0,-3,0,4),\qquad e_2=(0,0,1,0),
\end{equation}
and the real plus tensor
\begin{equation}
\varepsilon_+=\frac1{25}e_1e_1-e_2e_2.
\end{equation}
These are the convention-fixed polarizations entering \eqref{eq:AK}.

\subsection{Gauge and adjoint checks}

For each exchange, the exact bordered solution obeys
\begin{equation}
CX=0,
\qquad
KX+C^T\lambda+J=0.
\end{equation}
The total amplitude with graviton gauge polarization is zero, while the
contact contribution alone is
\begin{equation}
-\frac{596335802837691}{361367187500000}\ne0.
\end{equation}
This cancellation is a sensitive check of relative signs and
combinatorics.  De Donder and axial constraints give identical physical
totals.

Under reversal $p_a\mapsto-p_a$, all contacts and cubic currents are
recomputed.  The action contains two- and four-derivative terms, and the
independently recomputed quadratic matrices obey $K(-P)=K(P)$.  Contact,
$s$, $t$, and $u$ pieces equal their forward values separately.  Since
the external linear tensors are real, this is the physical reverse
matrix element used in $T_{2,+}^{\dg}$.

\subsection{Pole regression}

An exact shell-preserving complex shift
\begin{equation}
p_1(z)=p_1+z\eta,\qquad k(z)=k-z\eta,
\qquad \eta=(3,0,4i,5),
\end{equation}
has $\eta^2=\eta\cdot p_1=\eta\cdot k=0$.  In the relevant channel
\begin{equation}
P(z)^2=\frac{25}{4}+15z
\end{equation}
hits the massive pole at $z_*=-7/20$.  The bordered kernel there has
exactly the five TT zero modes.  If $T$ collects them and $G_5$ is their
Gram matrix,
\begin{equation}
W=T^TK'(z_*)T=-\frac{15}{2}G_5
\end{equation}
entrywise, including 24 nontrivial identities.  The residue
\begin{equation}
-j_1^TW^{-1}j_2=-\frac{4551434384\,i}{328515625}\ne0
\end{equation}
equals the independently normalized cubic contraction divided by the
kinetic slope.  A direct near-pole numerical probe of the actual bordered
exchange agrees to relative error below $10^{-4}$.  This regression is
not used to set the convention of the compact Einstein-frame residue;
it validates the original-variable propagator machinery.

\section{Covariant grading and Fock-lift equations}
\label{app:grading}

G18a solves the commutant rather than assuming Schur's lemma numerically.
At a standard massive momentum, represent the spin-$2$ fiber by symmetric
traceless $3\times3$ tensors.  The three $SO(3)$ generators act as
\begin{equation}
J_i(T)=L_iT+TL_i^T.
\end{equation}
Solving $[X,J_i]=0$ for a general $5\times5$ matrix gives
$X=a_MI_5$.  On the two complex massless helicity fibers, the little-
group rotation generator is $\operatorname{diag}(2,-2)$, whose
commutant is diagonal.  Solving simultaneously with the mass Casimir on
the full seven-dimensional fiber gives
\begin{equation}
X=\operatorname{diag}(a_+,a_-,a_MI_5).
\end{equation}
Hermiticity and $X^2=I$ restrict the three scalars to signs.  The
helicity-exchange real structure sets $a_+=a_-$; the free signature fixes
$(a_+,a_-,a_M)=(1,1,-1)$.

For the Fock lift, cluster factorization is the monoid equation
\eqref{eq:monoid}.  Iterating it gives
\begin{equation}
j(n_h,n_M)=j(1,0)^{n_h}j(0,1)^{n_M}=(-1)^{n_M}.
\end{equation}
The exact script verifies multiplicativity for all partitions in a
finite table and the analytic monoid argument establishes it for all
particle numbers.

\section{BRST block and non-nullity}
\label{app:brst}

The minimal non-null calculation uses
\begin{equation}
J_{\min}=\begin{pmatrix}1&0\\0&-1\end{pmatrix},
\qquad
X=\begin{pmatrix}0&0\\i\AK&0\end{pmatrix}.
\end{equation}
Then
\begin{equation}
X^\sharp=J_{\min}X^{\dg}J_{\min}
=\begin{pmatrix}0&i\AK\\0&0\end{pmatrix},
\qquad
X^\sharp X=\begin{pmatrix}-\AK^2&0\\0&0\end{pmatrix}.
\end{equation}
For $\mathcal O=-2i\AK\sigma_x$,
\begin{equation}
J_{\min}\mathcal O^{\dg}J_{\min}=\mathcal O,
\qquad \mathcal O^2=-4\AK^2I.
\end{equation}

To encode the cohomology argument, take a canonical finite splitting into
two physical representatives and one contractible doublet $u\mapsto
v\mapsto0$.  In the ordered basis $(i,f,u,v)$,
\begin{equation}
Q=\begin{pmatrix}
0&0&0&0\\
0&0&0&0\\
0&0&0&0\\
0&0&1&0
\end{pmatrix},
\qquad Q^2=0.
\end{equation}
For a completely general $4\times4$ matrix $Y$, direct multiplication
shows
\begin{equation}
\bigl(QY+YQ\bigr)_{\rm phys,phys}=0.
\end{equation}
The upper-left physical block of the G17 obstruction is nonzero after
embedding, so it cannot be of this exact form.  This finite matrix is not
a replacement for gravitational BRST quantization; it records the
universal algebraic implication used after the Ward and reduced-helicity
checks have established that the amplitude descends to cohomology.

\section{Machine-verification index}
\label{app:checks}

All scripts named below are in \path{symbolic/} and print
\texttt{ALL PASS}.  Checks use SymPy exact arithmetic except for the
explicitly labelled near-pole regression.

\begin{center}
\footnotesize
\begin{tabular}{p{0.11\linewidth}p{0.50\linewidth}p{0.29\linewidth}}
\hline
Check & Claim & Artifact/type\\
\hline
G13a--b & two-point on-shell validation and cubic Ward identities &
\path{gravity_perturbiner.py}, exact\\
G13 & $Mhh=0$ for $5\times4$ physical polarization choices &
\path{verify_gravity_cubic.py}, exact\\
G14a--c & nonzero original-variable $MMM$, $MMh$, and reduced massive-pole
factorization & \path{verify_gravity_cubic.py}, exact\\
G14d--i & potential cancellation, compact $-\sqrt6/8$ and $-\sqrt2/8$
amplitudes, symbolic Ward identity, five-state sum, inverse kernel &
\path{verify_gravity_factorization.py}, exact\\
G15a--f & full physical contact+$s,t,u$, Ward, Bose, gauge, and threshold
checks & \path{verify_gravity_g15.py}, exact\\
G15g & shifted-pole kernel, Schur slope, residue; near-pole probe &
same script, exact plus numerical regression\\
G16 & crossing needed for the adjoint is included in G17; exhaustive
$250$-polarization scan is deferred & not a theorem dependency\\
G17a--d & exposed contact/exchanges, reverse process, EOM/Ward/gauge,
termwise quarter-turn & \path{verify_gravity_obstruction.py}, exact\\
G17e--f & exact shell obstruction and first-order metric-ambiguity
independence & same script, exact\\
G18a--b & one-particle commutant and cluster-factorizing Fock lift &
\path{verify_gravity_krein.py}, exact\\
G18c--e & non-null quadratic traces, charge no-go, and BRST block &
same script, exact\\
\hline
\end{tabular}
\end{center}

The analytic inputs not reducible to finite symbolic identities are
identified in the main text: the induced-representation interpretation
of the free commutant (Paper~IV), the consistent-truncation/LSZ argument,
real-analytic continuation on the physical shell, and the standard BRST
decoupling statement.  The symbolic checks test every convention-sensitive
coefficient consumed by those arguments.

\paragraph{Reproduction commands.}
From the repository root for this project:
\begin{verbatim}
python3 symbolic/verify_gravity_completion.py
python3 symbolic/verify_gravity_cubic.py
python3 symbolic/verify_gravity_factorization.py
python3 symbolic/verify_gravity_g15.py
python3 symbolic/verify_gravity_obstruction.py
python3 symbolic/verify_gravity_krein.py
\end{verbatim}
The full G15 script contains a slower symbolic shifted-pole step; G17
recomputes the exact physical certificate and the checks needed by the
master theorem in a compact import-safe four-point assembly.

\end{document}
